% ============================================================
%  प्रकाशन एवं अधिकार पृष्ठ (Imprint / Copyright page)
% ============================================================
\thispagestyle{empty}
\vspace*{\fill}
\noindent
{\sffamily\bfseries\large\color{aiblue} विज्ञान में कृत्रिम बुद्धिमत्ता}\\[2pt]
{\sffamily\color{aigrey} सिद्धांत, अनुप्रयोग एवं भविष्य}\\[1pt]
{\sffamily\footnotesize\color{aigrey} Artificial Intelligence in Science: Principles, Applications and Future}

\medskip
\noindent\rule{\linewidth}{0.4pt}

\medskip
\noindent
\textbf{लेखक (Author):} Krishna Kumar Godara\\
\textbf{संस्करण (Edition):} प्रथम संस्करण (First Edition), 2026\\
\textbf{लक्षित पाठक (Audience):} विद्यालयी छात्र एवं विज्ञान-प्रेमी विद्यार्थी\\
\textbf{ISBN:} 978-93-5300-000-4 \textit{(नमूना — आवंटन हेतु प्रस्तावित / sample, to be assigned)}

\medskip
\noindent
\textcopyright~2026 Krishna Kumar Godara. कुछ अधिकार सुरक्षित (Some rights reserved)।

\medskip
\noindent
यह पुस्तक \textbf{शैक्षिक एवं अव्यावसायिक उपयोग} (educational, non-commercial use)
हेतु \href{https://creativecommons.org/licenses/by-nc/4.0/}{Creative Commons
BY-NC 4.0} अनुज्ञप्ति (license) के अंतर्गत उपलब्ध है। उपयुक्त श्रेय (attribution)
के साथ इसका पुनः उपयोग एवं वितरण किया जा सकता है।

\medskip
\noindent
\textbf{अस्वीकरण (Disclaimer):} इस पुस्तक की सामग्री केवल शैक्षिक उद्देश्य से दी गई है।
यहाँ उल्लिखित उपकरण, वेबसाइट एवं सेवाएँ उनके सम्बंधित स्वामियों की हैं; उनका उपयोग
करते समय आयु एवं गोपनीयता सम्बन्धी नियमों का पालन करें।

\medskip
\noindent
\textbf{सम्पर्क (Contact):} \texttt{kkgodara2000@gmail.com}\\
\textbf{Portfolio:} \texttt{kkreboot.github.io/portfolio}

\medskip
\noindent
{\footnotesize\color{aigrey} XeLaTeX द्वारा संकलित (typeset with XeLaTeX) ·
मुख्य फ़ॉन्ट: Noto Serif Devanagari.}
\vspace*{\fill}
\cleardoublepage
